\section{Literature positioning}
\label{sec:literature-positioning}

This work positions itself at the intersection of logical decidability,
partial observability, and diagnosability, while deliberately departing
from algorithmic, probabilistic, and control-oriented traditions.
All positioning is descriptive and contrastive. References to external
literature serve exclusively to situate the present results and do not
import additional primitives, semantics, or methodological assumptions
into the OC/ETS framework.

\subsection{Decidability and indistinguishability}

Classical notions of decidability originate in logic and computability
theory, where decidability is defined relative to formal languages,
derivation systems, or algorithms. In contrast, the present work adopts
a purely \emph{logical} notion of decidability that is independent of
algorithms, complexity, or computability. Decidability is defined
solely as invariance of a predicate over observational equivalence
classes induced by admissible observations.

This framing is closest in spirit to indistinguishability arguments in
logic and model theory, where different models satisfy the same
observable conditions while differing on internal structure. However,
the present setting is neither syntactic nor language-based:
equivalence is induced by the ETS-defined observation operator $\Pi$
acting on continua, rather than by satisfaction of a theory or a set of
sentences.
\cite{hodges1993_model_theory,ebbinghaus_flum1995_finite_model_theory,nerode1958_linear_automaton_transformations}

\subsection{Partial observability in systems theory}

Partial observability has been extensively studied in systems theory,
control, and estimation, typically with the objective of reconstructing
or estimating internal states from observed outputs. Such approaches
rely on assumptions about system dynamics, noise models, observability
ranks, identifiability, or asymptotic convergence.

The present work explicitly rejects this framing. Incomplete
observability is treated as a structural condition of the OC/ETS
formalism, not as a technical limitation to be overcome. No notion of
state reconstruction, observer synthesis, filtering, convergence, or
identifiability is employed. Consequently, the results presented here
are orthogonal to classical observability theory: they delineate
logical limits that apply even in the absence of noise and even with
unbounded observation time.
\cite{milner1980_ccs,milner1989_communication_concurrency,van_glabbeek2001_ltbts1}

\subsection{Diagnosability and fault detection}

Diagnosability is most prominently developed in the context of discrete
event systems, where it typically refers to the algorithmic ability to
infer fault occurrence from observed event sequences within bounded
delay.

The STOP-only results of this paper are superficially reminiscent of
diagnosability results in that they identify a class of terminal or
boundary conditions that can be certified from observations. The
similarity, however, is purely structural. No automata, formal
languages, delays, detection horizons, or decision algorithms are
involved, and no notion of detection time is defined. STOP-only
decidability is a logical classification result: it asserts invariance
of STOP across observational equivalence classes, not the existence of
a detection or decision procedure.
\cite{sampath1995_diagnosability_des}

\subsection{Epistemic limits and negative results}

The emphasis on negative and limit results places this work within a
broader tradition of epistemic limit analysis in logic and complex
systems. Rather than proposing improved methods, richer data, or
refined models, the paper identifies classes of questions that are
provably undecidable within a fixed formal system.

This positioning is intentional. Undecidability is not interpreted as
practical difficulty, lack of data, or methodological immaturity, but
as a formal impossibility statement relative to OC Core and
ETS.K\_EA.v1.1. No probabilistic, approximate, or asymptotic weakening
of these limits is admitted.
\cite{hodges1993_model_theory}

\subsection{Relation to enterprise modeling}

Within the enterprise modeling literature, most approaches focus on
representation, optimization, governance, or decision support. Even
when formal models are employed, they are typically embedded in
prescriptive or managerial narratives.

The present work adopts a different stance. Enterprises are treated as
continua equipped with admissible observability, and the sole question
addressed is what can be \emph{decided} about their states from
admissible observations. No claims are made about how enterprises
should be managed, intervened upon, optimized, or controlled. In this
sense, the contribution is foundational rather than applied.

\subsection{Bibliographic status}

This section is intentionally non-exhaustive. All references are
restricted to formally relevant literature on logical
indistinguishability, diagnosability, and epistemic limits, and are
included solely for positioning purposes. No argument in this paper
depends on results external to OC Core and ETS.K\_EA.v1.1.
